\section{Program-Verify Loops and Write Latency}
\label{sec:bg-pvloops}

Because OxRAM switching is stochastic, resulting in a wide variety of CF
states, a single programming pulse does not reliably bring a cell to a
target resistance state. Commercial architectures employ a
program-verify (PV) loop: a sequence of programming pulses interleaved
with resistance measurements that continues until the cell resistance
falls within a target window~\cite{
	kawahara2013filament,hayakawa2015highly,grossi2015impact,mao2016optimizing,
	milo2021accurate, hung2021mlc}. Each iteration of the loop applies a
pulse and checks the result. If the cell has not reached the target
state, another pulse is applied. The loop terminates when the cell
resistance is confirmed to be within specification. Single-level cells
(SLCs) define two resistance windows corresponding to logical 0 and 1,
while multi-level cells (MLCs) require finer resistance windows and more
sophisticated verification schemes but allow for the encoding of
multiple bits per cell~\cite{hung2021mlc}.

The non-deterministic completion time of a PV loop contributes to the
write latency $t_w$ of commercial ReRAM memories. This stochastic timing
contribution arises from cycle-to-cycle variability in filament
switching dynamics. This variabilily effects the number of pulses
required to grow or rupture a filament varies even under identical
initial conditions ~\cite{yu2012switching, kawahara2013filament}. PV
loops typically apply different different biases to each cell depending
on the operation being performed. In 1T1R ReRAM arrays (described in the
next section), these biases are chosen such that access transistors can
serve as current compliance to avoid overshoots during SET operations.
At the word-level, a single write operation may require both SETs and
RESETs. In this case, the write controller will must serialize SET/RESET
pulses at the cost of increased latency. Despite some simulated results
that suggest word-level biases can be achieved without this
serialization \cite{mao2016optimizing}, ReRAM studies using physical
ReRAM typically calibrate SET and RESET pulses
separately~\cite{hayakawa2017resolving}.

The data-dependence of write latency has direct implications for timing
side-channel analysis. Writes that require only SET transitions, only
RESET transitions, a combination of both, or no transitions at all
produce different latency distributions. In a crossbar architecture,
writes that require both SET and RESET transitions within the same word
incur additional serialization overhead, as described in
Section~\ref{sec:bg-crossbar}. The presence of an on-die ECC further
modifies the effective transition composition of every write, since
parity bits computed over user data may require transitions that differ
in polarity from the data transitions. The net result is that the write
latency of a commercial ReRAM device encodes information about the
codeword being written, not merely the user-visible data word. This
relationship is the basis of the BREW-RC technique presented in
Chapter~\ref{chap:brewrc}.